% Section 6: the costs, and where the two algorithms sit in the literature.
% Sources are marked peer-reviewed or practitioner, because the maze-generation
% literature is unusually thin and the distinction matters.
\section{Complexity, and what the literature says}
\label{sec:analysis}

\subsection{Size of the data structure}

Write $n=WH$ for the number of cells.

\begin{center}
\begin{tabular}{@{}lll@{}}
  \toprule
  What & How many & Why \\
  \midrule
  Nodes                  & $n$        & one per cell \\
  Neighbour entries      & $2(n-1)$   & $n-1$ doors, each recorded at both ends \\
  Index $\langle P_i,P_j\rangle$ & $n$ & one reference per cell \\
  Region tree nodes      & $2L-1$     & $L$ corridor leaves, $L-1$ divisions \\
  \bottomrule
\end{tabular}
\end{center}

Everything is $\Theta(n)$. The grid graph is sparse --- $|\gridedges|=2n-W-H$
--- so no representation of it can be smaller than linear, and the maze itself
is sparser still.

\subsection{Cost of generation}

Each call of \texttt{Divide} does $O(1)$ work of its own: two draws, one
\texttt{Link}, and the recursion. Each call of \texttt{SpanCorridor} does work
proportional to the length of its corridor.

Every door is opened exactly once, by one division or by one corridor
(Section~\ref{sec:one-door-per-division}), and there are $n-1$ of them. The
region tree has $2L-1$ nodes with $L\le n$. So

\begin{keypoint}
Generation runs in \keyterm{$\Theta(n)=\Theta(WH)$ time}, which is optimal: any
generator must write $n-1$ doors.
\end{keypoint}

The recursion is not balanced. The wall position is uniform, not central, so a
division can leave a single column on one side --- the same skew that makes
quicksort's worst case quadratic. Depth is $O(n)$ in the worst case and
$O(\log n)$ in expectation, and the constant is visible:

\begin{center}
\begin{tabular}{@{}lrrr@{}}
  \toprule
  Area & Corridor leaves & Mean depth & Depth if balanced \\
  \midrule
  $16\times16$   & $55.2$   & $11.6$ & $5.8$ \\
  $64\times64$   & $671.7$  & $21.2$ & $9.4$ \\
  $128\times128$ & $2388.5$ & $26.7$ & $11.2$ \\
  \bottomrule
\end{tabular}
\end{center}

Roughly $2.4$ times the depth of a balanced tree on the same number of leaves,
at every size measured. That is the cost of drawing the wall uniformly, and it
is the figure a stack-depth bound has to be written against.

Stopping at corridors changes no asymptotics and a great deal of constant:
$1301$ calls instead of $8191$ on $64\times64$, and one call instead of $79$ on
$1\times40$.

\subsection{Cost of solving}

The maze is a tree, so $|E|=n-1$ and the graph the search runs on is as sparse
as a connected graph can be. With a binary heap, Dijkstra costs
$O((|V|+|E|)\log|V|)=O(n\log n)$~\cite[Ch.~22]{cormen2022introduction}.

That $\log$ buys nothing here. All weights are $1$, so Dijkstra settles cells in
depth order and is breadth-first search wearing a priority queue; plain
breadth-first search would do the same work in $\Theta(n)$. The queue is kept
because the report's algorithm is written to accept an estimate, and A* needs
it.

\subsection{Why the heuristic barely helps}

The measurements of Section~\ref{sec:cpp-testing} are not a disappointing
implementation. They are what the search literature predicts for mazes, for two
separate reasons that are worth keeping apart.

\begin{description}[leftmargin=*,style=nextline,itemsep=6pt]
  \item[A tie-break cannot prune.] Every cell nearer the start than the exit is
    \emph{necessarily expanded} by any admissible search with this priority
    function~\cite{dechter1985generalized}, so only the last layer is
    negotiable. Proposition~\ref{prop:tiebreak-ceiling} is that result applied
    here, and the measured $0.03\%$ is it.
  \item[A* prunes subtrees, not rival routes.] On a general grid much of A*'s
    advantage is in discarding the many equal-length routes between the same
    pair of cells. A perfect maze has \emph{none}: it is a tree, so there are no
    rival routes to discard. What is left is pruning dead-end branches whose
    depth plus estimate already exceeds $d$ --- real, but a smaller effect.
  \item[The estimate is very inaccurate on mazes.] This is the measured result
    in the literature, not an intuition. Sturtevant et
    al.~\cite{sturtevant2009memory} report that on $512\times512$ maze maps the
    geometric estimate averages $151$ where the true distance is $512$--$768$,
    and that a memory-based heuristic expands over eleven times fewer nodes. A
    straight line between two cells says little about a route that must wind.
\end{description}

\begin{keypoint}
The walls, not the algorithm, are what defeats the guide. The estimate ignores
them, and in a maze they are the whole problem.
\end{keypoint}

A consequence worth stating for completeness: the grid-pathfinding techniques
that give the largest speedups --- Jump Point Search and the symmetry-breaking
family~\cite{harabor2011online} --- work by collapsing equal-cost alternative
routes. A perfect maze has no path symmetry at all, so that line of work has
nothing to prune here. Standard benchmark sets include maze maps for exactly
this reason~\cite{sturtevant2012benchmarks}.

\subsection{Where recursive division sits}

The maze-generation literature is small, and mostly practitioner-written; the
distinction is marked below because it matters for how much weight each claim
carries.

\paragraph{The vocabulary.}
What this report builds is a \keyterm{perfect maze} --- a maze in which any two
cells are joined by a unique path --- which is the same object as a spanning
tree of the grid graph. The term is Pullen's~\cite{pullen2026think} and is used
verbatim in peer-reviewed work~\cite{bellot2021how}. Recursive division is a
\keyterm{wall adder}: it places walls, where most generators carve passages.
Bellot et al.\ call it the only wall adder among the standard algorithms.

\paragraph{The alternatives.}
Randomised depth-first search, randomised Kruskal and Prim, Eller's algorithm,
Hunt-and-Kill, binary tree and sidewinder all produce spanning trees of the same
grid, in $\Theta(n)$ or close to it. Two do something this one does not:
Aldous--Broder~\cite{aldous1990random,broder1989generating} and
Wilson's algorithm~\cite{wilson1996generating} sample a \keyterm{uniform}
spanning tree --- every perfect maze equally likely. Wilson's runs in expected
mean-hitting time rather than cover time, which on a grid is a $\Theta(\log n)$
improvement over Aldous--Broder.

\paragraph{Recursive division is not uniform, and not nearly.}
This is the honest position of the algorithm chosen here. Its first division
leaves exactly one door crossing a full straight line, and recursively so; any
spanning tree admitting no such decomposition can \emph{never} be produced. The
support is a strict subset of the spanning trees, so the distribution is not
uniform even before asking how it is weighted within its support.
Pullen~\cite{pullen2026think} classifies it accordingly, and records its visible
signature: long straight walls crossing the interior, rectangular rooms, a
self-similar texture, and a notably direct solution path. The count makes the
gap concrete --- an $n$-cell grid has on the order of $3.21^{\,n}$ spanning
trees, and recursive division reaches only those with a guillotine
decomposition.

\begin{keypoint}
Recursive division is chosen for being \keyterm{simple, linear and recursive},
which is what this APP is about. It is not chosen for the quality of its
randomness, and it should not be claimed for it.
\end{keypoint}

\paragraph{What is missing from the literature.}
No peer-reviewed analysis of the distribution recursive division induces appears
to exist. Bellot et al.~\cite{bellot2021how} include it but score only
difficulty; Gabrov\v{s}ek~\cite{gabrovsek2019analysis} measures six generators
and omits it. The characterisation of its bias is practitioner-level throughout.
Measuring its degree distribution and straight-run lengths against the uniform
spanning tree would close that gap, and is the natural continuation of this
work.
